\chapter{Iteración 7}

\begin{quotation}[Empresario]{Steve Jobs (1955--2011)}
El diseño no es solo lo que ves, sino cómo funciona.
\end{quotation}

\begin{abstract}
En este capítulo reflejaremos todos y cada uno de los costes que se han de reflejar a la hora de desarrollar el proyecto así como la relación de funcionalidades que se desarrollan con estos costes.
\end{abstract}

\section{Características a desarrollar}

\begin{enumerate}
\item Funcionalidad 1: Implementación del dashboard. Ver Tabla \ref{tab:valor1AportadoIteracion7}.
\end{enumerate}
%----------------------------------------------------------------------------------------------------------
\begin{table*}[htb]
	\centering
	\begin{coolTable}{p{4cm}p{\textwidth-4.5cm}}{2}
{Análisis de valor aportado 0001}
\textbf{Propuesta}&Se pretende desarrollar un dashboard con los datos relevantes del usuario, así como los avances diarios en cuanto a la realización de tareas.\\
\midrule
\textbf{Valor}&El valor aportado es poder visualizar los datos relevantes del usuario, así como los avances diarios en cuanto a la realización de tareas.\\
\textbf{Coste}&Se requiere la participación de ambos miembros del equipo desarrollador para llevar a cabo la implementación.\\
\textbf{Opciones}&Esta funcionalidad tiene un valor alto en la aplicación.\\
\textbf{Riesgos}&La no correcta visualización de los diferentes campos con los que cuenta el dashboard.\\
\textbf{Deuda técnica}&Se ha elegido la solución más factible para el desarrollo de la misma.\\
	\end{coolTable}
	\caption{Análisis de valor aportado 0001\label{tab:valor1AportadoIteracion7}}
\end{table*}
%----------------------------------------------------------------------------------------------------------

\clearpage

\section{Diseño}

Esto es comparable en la siguiente figura. \ref{fig:diseno07}.

\begin{figure}[htbp]
\begin{center}

\includegraphics[width=\textwidth]{figures/UML-iteracion-6.jpg}
\caption{Diagrama UML de diseño para la iteración 7}
\label{fig:diseno07}
\end{center}
\end{figure}

\clearpage


\begin{table*}[htb]
	\centering
	\begin{coolTable}{p{4cm}p{\textwidth-4.5cm}}{2}
{Memorando técnico 0001}
\textbf{Asunto}&Diseño del modelo Tarea.\\
\textbf{Resumen}&Se ha diseñado un modelo de Tarea que contenga los elementos necesarios para nuestra definición de esta; título, descripción, fechas de inicio y fin, duración, priorización si está terminada o no, la estimación de la finalización de la tarea, además contaremos con el creador de la tarea, el usuario asignado a la tarea y el equipo de la tarea si esta tuviera.\\
\midrule
\textbf{Factores causantes}&La relación entre modelos, ya que la clase Tarea cuenta con una asociación tanto con el modelo Equipo como con el modelo Usuario, así como los datos que requerimos que estos tengan asignados para el proyecto. \\
\textbf{Solución}&El modelo resultante asigna un tipo de dato para cada uno de los campos así como las relaciones que no pueden ser nulas en ningún momento por la concepción del modelo de Tarea como pueden ser el nombre o el creador.\\
\textbf{Motivación}&Un modelo que contenga los datos necesarios para un correcto procesamiento de las tareas.\\
\textbf{Cuestiones abiertas}& La correcta relación entre tablas y completitud de los atributos.\\
\textbf{Alternativas}&Se valoraron otras alternativas de relaciones, pero la dispuesta es la que se consideró más óptima tanto a niveles de implementación como de seguridad.\\
	\end{coolTable}
	\caption{Memorando técnico 0001}
\end{table*}

%-----------------------------------------------------------------------------------------------------------
\begin{table*}[htb]
	\centering
	\begin{coolTable}{p{4cm}p{\textwidth-4.5cm}}{2}
{Memorando técnico 0002}
\textbf{Asunto}&Diseño del modelo Equipo.\\
\textbf{Resumen}&Se ha diseñado un modelo de Equipo que contenga los elementos necesarios para nuestra definición de esta; nombre, además contaremos con el creador del equipo, lista de asociados, así como de la lista de managers.\\
\midrule
\textbf{Factores causantes}&La relación entre modelos, ya que la clase Equipo cuenta con una asociación con el modelo Usuario, así como los datos que requerimos que estos tengan asignados para el proyecto.\\
\textbf{Solución}&El modelo resultante asigna un tipo de dato para cada uno de los campos.\\
\textbf{Motivación}&Un modelo que contenga los datos necesarios para un correcto procesamiento de los equipos.\\
\textbf{Cuestiones abiertas}& La correcta relación entre tablas y completitud de los atributos.\\
\textbf{Alternativas}&Dado a la tecnología utilizada la manera de llevar a cabo esta relación entre tablas es la más óptima.\\
	\end{coolTable}
	\caption{Memorando técnico 0002}
\end{table*}

%-------------------------------------------------------------------------------------------------------

\begin{table*}[htb]
	\centering
	\begin{coolTable}{p{4cm}p{\textwidth-4.5cm}}{2}
{Memorando técnico 0003}
\textbf{Asunto}&Diseño del modelo Usuario.\\
\textbf{Resumen}&Se ha diseñado un modelo de Usuario que contenga los elementos necesarios para nuestra definición de este; nombre, apellidos, correo electrónico, la hora de inicio de su jornada laboral, hora de fin de la jornada laboral, nombre de usuario, contraseña. Se ha decidido remover el atributo boolean para comprobar si se había iniciado sesión o no debido a que el token \emph{Bearer}, de la misma manera ya que no se requería el atributo de los equipos porque con la propia asociación se obtienen esos datos siendo innecesario el atributo.\\
\midrule
\textbf{Factores causantes}&La relación entre modelos, ya que la clase Usuario cuenta con dos asociaciones con el modelo Tarea y con el modelo Equipo, así como los datos que requerimos que estos tengan asignados para el proyecto.\\
\textbf{Solución}&El modelo resultante asigna un tipo de dato para cada uno de los campos necesarios resultando en la eliminación de estos dos atributos, logged y Equipos.\\
\textbf{Motivación}&Un modelo que contenga los datos necesarios para un correcto procesamiento de las Usuario.\\
\textbf{Cuestiones abiertas}&La correcta relación entre tablas y completitud de los atributos.\\
\textbf{Alternativas}&Se valoraron otras alternativas de relaciones, pero la dispuesta es la que se consideró más óptima tanto a niveles de implementación como de seguridad.\\
	\end{coolTable}
	\caption{Memorando técnico 0003}
\end{table*}

%----------------------------------------------------------------------------------------------------------
\begin{table*}[htb]
	\centering
	\begin{coolTable}{p{4cm}p{\textwidth-4.5cm}}{2}
{Memorando técnico 0004}
\textbf{Asunto}&Diseño del modelo Invitación.\\
\textbf{Resumen}&Se ha diseñado un modelo de Invitación que contenga los elementos necesarios para nuestra definición de esta; mensaje, estado, además contaremos con el usuario invitado, así como el equipo al que pertenece la invitación.\\
\midrule
\textbf{Factores causantes}&La relación entre modelos, ya que la clase Invitación cuenta con una asociación con el modelo Usuario y con una asociación con el modelo Equipo, así como los datos que requerimos que estos tengan asignados para el proyecto.\\
\textbf{Solución}&El modelo resultante asigna un tipo de dato para cada uno de los campos.\\
\textbf{Motivación}&Un modelo que contenga los datos necesarios para un correcto procesamiento de las invitaciones.\\
\textbf{Cuestiones abiertas}& La correcta relación entre tablas y completitud de los atributos.\\
\textbf{Alternativas}&Dado a la tecnología utilizada la manera de llevar a cabo esta relación entre tablas es la más óptima.\\
	\end{coolTable}
	\caption{Memorando técnico 0004}
\end{table*}

%------------------------------------------------------------------------------------------------------
\begin{table*}[htb]
	\centering
	\begin{coolTable}{p{4cm}p{\textwidth-4.5cm}}{2}
{Memorando técnico 0005}
\textbf{Asunto}&Diseño del modelo Notificación.\\
\textbf{Resumen}&Se ha diseñado un modelo de Notificación que contenga los elementos necesarios para nuestra definición de esta; mensaje, estado, tipo de notificación, además contaremos con el usuario invitado, así como el equipo al que pertenece la invitación.\\
\midrule
\textbf{Factores causantes}&La relación entre modelos, ya que la clase Notificación cuenta con una asociación con el modelo Usuario, así como los datos que requerimos que estos tengan asignados para el proyecto.\\
\textbf{Solución}&El modelo resultante asigna un tipo de dato para cada uno de los campos.\\
\textbf{Motivación}&Un modelo que contenga los datos necesarios para un correcto procesamiento de las notificaciones.\\
\textbf{Cuestiones abiertas}& La correcta relación entre tablas y completitud de los atributos.\\
\textbf{Alternativas}&Dado a la tecnología utilizada la manera de llevar a cabo esta relación entre tablas es la más óptima.\\
	\end{coolTable}
	\caption{Memorando técnico 0005}
\end{table*}

%------------------------------------------------------------------------------------------------------

\clearpage

\section{Implementación}

En esta sección se va a realizar un memorando técnico por cada decisión de implementación y refactorización que afecte al diseño.

\begin{table*}[htb]
	\centering
	\begin{coolTable}{p{4cm}p{\textwidth-4.5cm}}{2}
{Memorando técnico 0001}
\textbf{Asunto}&La creación de la clase Tarea.\\
\textbf{Resumen}&Se han definido los siguientes atributos para la clase Tarea que son los necesarios para mostrar los datos que queremos.\\
\midrule
\textbf{Factores causantes}&La relación entre tablas, ya que la clase Tarea cuenta con una asociación tanto con la clase Equipo como con la clase Usuario. \\
\textbf{Solución}&La solución llevada a cabo tiene la utilización de las anotaciones que proporciona JPA, para gestionar las relaciones entre dos tablas. En este caso para la relación entre Tarea y Equipo será \textbf{@ManyToOne} y para la relación entre Tarea y Usuario será \textbf{@OneToOne}.\\
\textbf{Motivación}&Ya conocíamos esta forma de relacionar las tablas.\\
\textbf{Cuestiones abiertas}& La correcta relación entre tablas.\\
\textbf{Alternativas}&Dado a la tecnología utilizada la manera de llevar a cabo esta relación entre tablas es la más óptima.\\
	\end{coolTable}
	\caption{Memorando técnico 0001}
\end{table*}

%-----------------------------------------------------------------------------------------------------------
\begin{table*}[htb]
	\centering
	\begin{coolTable}{p{4cm}p{\textwidth-4.5cm}}{2}
{Memorando técnico 0002}
\textbf{Asunto}&La creación de la clase Equipo.\\
\textbf{Resumen}&Se han definido los siguientes atributos para la clase Equipo que son los necesarios para mostrar los datos que queremos.\\
\midrule
\textbf{Factores causantes}&La relación entre tablas, ya que la clase Equipo cuenta con una asociación con la clase Usuario. \\
\textbf{Solución}&La solución llevada a cabo tiene la utilización de las anotaciones que proporciona JPA, para gestionar las relaciones entre dos tablas. En este caso la relación entre Equipo y Usuario serán \textbf{@ManyToMany}.\\
\textbf{Motivación}&Ya conocíamos esta forma de relacionar las tablas.\\
\textbf{Cuestiones abiertas}& La correcta relación entre tablas.\\
\textbf{Alternativas}&Dado a la tecnología utilizada la manera de llevar a cabo esta relación entre tablas es la más óptima.\\
	\end{coolTable}
	\caption{Memorando técnico 0002}
\end{table*}

%-------------------------------------------------------------------------------------------------------

\begin{table*}[htb]
	\centering
	\begin{coolTable}{p{4cm}p{\textwidth-4.5cm}}{2}
{Memorando técnico 0003}
\textbf{Asunto}&La creación de la clase Usuario.\\
\textbf{Resumen}&Se han definido los siguientes atributos para la clase Usuario que son los necesarios para mostrar los datos que queremos.\\
\midrule
\textbf{Factores causantes}&La relación entre tablas, ya que la clase Usuario cuenta con una asociación con la clase Equipo. \\
\textbf{Solución}&La solución llevada a cabo tiene la utilización de las anotaciones que proporciona JPA, para gestionar las relaciones entre dos tablas. En este caso para la relación entre Usuario y Equipo será \textbf{@OneToMany}.\\
\textbf{Motivación}&Ya conocíamos esta forma de relacionar las tablas.\\
\textbf{Cuestiones abiertas}& La correcta relación entre tablas.\\
\textbf{Alternativas}&Dado a la tecnología utilizada la manera de llevar a cabo esta relación entre tablas es la más óptima.\\
	\end{coolTable}
	\caption{Memorando técnico 0003}
\end{table*}

%----------------------------------------------------------------------------------------------------------
\begin{table*}[htb]
	\centering
	\begin{coolTable}{p{4cm}p{\textwidth-4.5cm}}{2}
{Memorando técnico 0004}
\textbf{Asunto}&La creación de la clase Invitación.\\
\textbf{Resumen}&Se han definido los siguientes atributos para la clase Invitación que son los necesarios para mostrar los datos que queremos.\\
\midrule
\textbf{Factores causantes}&La relación entre tablas, ya que la clase Invitación cuenta con una asociación tanto con la clase Equipo como con la clase Usuario. \\
\textbf{Solución}&La solución llevada a cabo tiene la utilización de las anotaciones que proporciona JPA, para gestionar las relaciones entre dos tablas. En este caso para la relación entre Invitación y Equipo será \textbf{@OneToOne} y para la relación entre Invitación y Usuario será \textbf{@OneToOne}.\\
\textbf{Motivación}&Ya conocíamos esta forma de relacionar las tablas.\\
\textbf{Cuestiones abiertas}& La correcta relación entre tablas.\\
\textbf{Alternativas}&Dado a la tecnología utilizada la manera de llevar a cabo esta relación entre tablas es la más óptima.\\
	\end{coolTable}
	\caption{Memorando técnico 0004}
\end{table*}

%-----------------------------------------------------------------------------------------------------
\begin{table*}[htb]
	\centering
	\begin{coolTable}{p{4cm}p{\textwidth-4.5cm}}{2}
{Memorando técnico 0005}
\textbf{Asunto}&La creación de la clase Notificación.\\
\textbf{Resumen}&Se han definido los siguientes atributos para la clase Invitación que son los necesarios para mostrar los datos que queremos.\\
\midrule
\textbf{Factores causantes}&La relación entre tablas, ya que la clase Notificación cuenta con una asociación con la clase Usuario. \\
\textbf{Solución}&La solución llevada a cabo tiene la utilización de las anotaciones que proporciona JPA, para gestionar las relaciones entre dos tablas. En este caso para la relación entre Notificación y Usuario será \textbf{@ManyToOne}.\\
\textbf{Motivación}&Ya conocíamos esta forma de relacionar las tablas.\\
\textbf{Cuestiones abiertas}& La correcta relación entre tablas.\\
\textbf{Alternativas}&Dado a la tecnología utilizada la manera de llevar a cabo esta relación entre tablas es la más óptima.\\
	\end{coolTable}
	\caption{Memorando técnico 0005}
\end{table*}

%-----------------------------------------------------------------------------------------------------

\clearpage

\section{Pruebas}

La totalidad de las pruebas se han realizado en la iteración 6, en esta iteración no se han desarrollado nuevos métodos que impliquen el desarrollo de pruebas puesto que todas las funcionalidades a desarrollar pertenecían al Front-End.

\section{Despliegue}

Debido a que se ha producido casi la completitud del código en general y en particular durante la 6ª iteración se ha completado el desarrollo de todo el código en Back-End durante la primera semana de esta iteración no se ha desarrollado ningún despliegue. Sin embargo, durante la segunda semana el lunes 29 de agosto se ha desplegado la versión definitiva del Back-End en la siguiente url: \url{https://t-planifica.herokuapp.com/}.

